Our goal in this section is to prove the following theorem.

\begin{theorem}{Existence and uniqueness of weak solutions}{eq_and_un}
    Let the initial conditions $u_0, u_{b, 0} \in \Ltwo{\partial \Omega}$ and $v_0 \in \Ltwo{\Omega}$ be essentially bounded and non-negative. Then there exist unique weak solutions $u, u_b \in H^1(0, T; \Ltwo{\Gamma})$ and $v \in \Ltwo{0, T; H^1(\Omega)}$ with $v_t \in \Ltwo{0, T; H^1(\Omega)^*}$ to \cref{sys:ODEs,sys:PDEs,sys:Init}. Moreover, these solutions are also non-negative and bounded.
\end{theorem}

As previously mentioned, the proof will be based on an application of Banach's fixed-point theorem. The proof for our model follows the same general idea as the proof of the Picard-Lindel\"of theorem proving the existence and uniqueness for ODEs, but we need to adjust the arguments and function spaces to handle our situation of a coupled bulk-surface system. The idea is to consider the problem at a smaller time interval $[0, \widetilde{T}]$ for some $\widetilde{T} \leq T$ and ``freeze'' the function $v$ in the ODEs so we can use the standard existence and uniqueness theory for ODEs to get solutions $u$ and $u_b$ to this problem. The PDE we treat similarly, freezing the functions $u$ and $u_b$. Having a solution of this PDE, we can ``update'' the frozen terms in the ODE. This produces an operator $\mathcal{T}$ that maps any $v \in \pdespace$ to another $\widetilde{v} \in \pdespace$, where $\pdespace$ is the space defined by
%
\begin{equation*}
    \pdespace := \left\{ v \in \Ltwo{0, \widetilde{T}; H^1(\Omega)}\ \middle|\ v_t \in \Ltwo{0, \widetilde{T}; H^1(\Omega)^*} \text{ and } 0 \leq v \leq M_v(T) \text{ almost everywhere} \right\},
\end{equation*}

\noindent
which we view as a closed subset of $\Ltwo{0, \widetilde{T}; \Ltwo{\Omega}}$. We will derive the exact value of $M_v(T)$ later. Defining this operator, we will need another function space $\odespace$ given by
%
\begin{equation*}
    \odespace := \left\{ (u, u_b) \in \left(H^1(0, \widetilde{T}; \Ltwo{\Gamma})\right)^2\ \middle|\ 0 \leq u \leq M_u(T) \text{ and } 0 \leq u_b \leq M_{u_b}(T) \text{ almost everywhere} \right\},
\end{equation*}

\noindent
which we view as a closed subset of the space $(\Ltwo{0, \widetilde{T}; \Ltwo{\Gamma}})^2$. Again, the exact values of $M_u(T)$ and $M_{u_b}(T)$ will be derived later. If we can then show that this operator is a contraction, Banach's fixed-point theorem will guarantee the existence and uniqueness of a fixed point $v^* \in \pdespace$, which will be a local solution to the PDE. We will then show how we can stitch together local solutions to get a unique solution for every final time $T$.

Let us first make precise what we mean by ``freezing'' the function $v$ in the ODE system. Fix $\widetilde{T} \leq T$ and $v^* \in \pdespace$ and consider the modified system of ODEs given by
%
\begin{equation}\label{eq:frozen_odes}
    \begin{aligned}
        u_t         &= P_u + \alpha u_b - \delta_u u - 2 \kon u^2 \gamma v^* + 2 \koff u_b   & \text{on }& (0, \widetilde{T}) \times \Gamma, \\
        u_{b,t}     &= - \delta_{u_b} u_b + \kon u^2 \gamma v^* - \koff u_b                  & \text{on }& (0, \widetilde{T}) \times \Gamma, \\
        u(0)        &= u_0, \qquad u_b(0) = u_{b, 0}                                  & \text{on }& \Gamma. \\
    \end{aligned}
\end{equation}

\noindent
These ODEs are now \emph{decoupled} from the PDE, so we can investigate the existence and uniqueness of solutions to this system without simultaneously having to prove that $v$ is a solution of the PDE.

\begin{lemma}{Existence and uniqueness for the decoupled ODEs}{exun_frozen_odes}
    Let the initial conditions $u_0, u_{b,0} \in \Ltwo{\Gamma}$ be essentially bounded and non-negative and fix $v^* \in \pdespace$. Then there exist unique weak solutions $(u, u_b) \in \odespace$ to \cref{eq:frozen_odes} satisfying these initial conditions.
\end{lemma}

\begingroup
\newcommand{\vb}{\overline{v}}
\newcommand{\vt}{\widetilde{v}}
\newcommand{\ub}{\overline{u}}
\newcommand{\ut}{\widetilde{u}}
\newcommand{\ubb}{\overline{u}_b}
\newcommand{\ubt}{\widetilde{u}_b}
\begin{proof}
    Fix $x$ and define $w$ and $F: \R \times \R^2 \to \R^2$ as in \cref{eq:weak_ode,eq:ode_system_function}, where $v$ is replaced by $v^*$. Since the components of $F$ are polynomial in $w$, the function $F$ is certainly continuous in $w$. Additionally, by the regularity of $\gamma v^*$ which we showed in \cref{prop:trace_regularity}, we know that $F$ is measurable in $t$. This also implies that for any compact subset of $\R \times \R^2$, we can bound $F$ by some time-dependent integrable function $m$. In particular, this subset can be taken to be a rectangle\footnote{On first sight, it can seem strange to call this type of set a rectangle instead of a cylinder. This name comes from the fact that in our main source \cite{coddington_theory_1984}, the theorem is proved for the case of $n = 1$. } as in the hypothesis of \cref{thm:caratheodory}, so let $R$ be such a rectangle around the initial time $t = 0$ and initial data $(u_0, u_{b,0})$. The hypotheses of \cref{thm:caratheodory} are then fully satisfied, giving us a local solution $w: [0, \beta] \to \R^2$ for some $\beta > 0$, with $w(0) = (u_0, u_{b,0})$.

    We show that this local solution is unique. Suppose we have two such local solutions, $w_1$ and $w_2$. Since these solutions are absolutely continuous, they are bounded on $[0, \beta]$, so we can take a compact set $K \subseteq \R^2$ that contains both their trajectories. For any $t \in [0, \beta]$ and $\overline{w} = (\ub, \ubb), \widetilde{w} = (\ut, \ubt) \in K$, $F$ satisfies
%
    \begin{equation*}
    \begin{aligned}
        \left|F(t, \overline{w}) - F(t, \widetilde{w})\right|^2 &= \left| \begin{pmatrix}
            -2 \kon \left(\ub^2 - \ut^2\right) \gamma v^*(t, x) + 2 \koff \left(\ubb - \ubt\right) + \alpha \left(\ubb - \ubt\right) - \delta_u \left(\ub - \ut\right) \\
            \kon \left(\ub^2 - \ut^2\right) \gamma v^*(t, x) - \koff \left(\ubb - \ubt\right) - \delta_{u_b} \left(\ubb - \ubt\right) \\
        \end{pmatrix} \right|^2 \\
        &= \left| \begin{pmatrix}
            - \left(2 \kon \left(\ub + \ut\right) \gamma v^*(t, x) + \delta_u\right) \left(\ub - \ut\right) & + \left(2 \koff + \alpha\right)\left(\ubb - \ubt\right) \\
            \kon \left(\ub + \ut\right) \gamma v^*(t, x) \left(\ub - \ut\right) & - \left(\koff + \delta_{u_b}\right) \left(\ubb - \ubt\right) \\
        \end{pmatrix} \right|^2 \\
        &\leq 2 \left| \begin{pmatrix}
            - 2 \kon \left(\ub + \ut\right) \gamma v^*(t, x) - \delta_u \\
            \kon \left(\ub + \ut\right) \gamma v^*(t, x)
        \end{pmatrix} \left(\ub - \ut\right) \right|^2 + 2 \left| \begin{pmatrix}
            2 \koff + \alpha \\
            \koff + \delta_{u_b}
        \end{pmatrix} \left(\ubb - \ubt\right) \right|^2 \\
        &\leq L_K(t)^2 \left(\left(\ub - \ut\right)^2 + \left(\ubb - \ubt\right)^2\right) \leq L_K(t)^2 \left|\overline{w} - \widetilde{w}\right|^2,
    \end{aligned}
    \end{equation*}

    \noindent
    where $L_K(t)$ is defined as
%
    \begin{equation*}
        L_K(t) := \sqrt{2} \max \left\{ \left| \begin{pmatrix}
            4 \kon \diam(K) v^*(t, x) + \delta_u \\
            2 \kon \diam(K) v^*(t, x)
        \end{pmatrix} \right| ,\quad \left| \begin{pmatrix}
            2 \koff + \alpha \\
            \koff + \delta_{u_b}
        \end{pmatrix} \right| \right\} ,
    \end{equation*}

    \noindent
    which is of $L^2$-regularity in $t$ due to the essential boundedness of $v^*$. Now for the difference of the two solutions we obtain
%
    \begin{equation*}
        \begin{aligned}
            \left|w_1(t) - w_2(t)\right| &= \left|(u_0, u_{b,0}) + \int_{0}^{t} F(t, w_1(t)) \d t - (u_0, u_{b,0}) - \int_{0}^{t} F(t, w_2(t)) \d t \right| \\
            &\leq \int_{0}^{t} \left|F(t, w_1(t)) - F(t, w_2(t))\right| \d t \leq \int_{0}^{t} L_K(t) \left|w_1(t) - w_2(t)\right| \d t
        \end{aligned}
    \end{equation*}

    \noindent
    for every $t \in [0, \beta]$. Using Gronwall's lemma and the fact that $w_1(0) = w_2(0)$, we see that this implies that $\left|w_1(t) - w_2(t)\right| = 0$ for all $t \in [0, \beta]$, so $w_1 = w_2$.

    To show that the components $u$ and $u_b$ of this solution $w$ are non-negative, we multiply the equations of \cref{eq:frozen_odes} by $u_{-} := - \min (0, u)$ and $u_{b, -} := - \min (0, u_b)$, respectively, and obtain
%
    \begin{align*}
        \left( u_{b,-}(t) \right)^2 &= 2 \int_{0}^{t} u_{b,-}(r) \frac{\d}{\d t} u_{b, -}(r)\d r = - 2 \int_{0}^{t} u_{b,-}(r) u_{b, t}(r) \d r \\
        &= - 2 \int_{0}^{t} \left[ - \delta_{u_b} u_b(r) + \kon u(r)^2 \gamma v^*(r, x) - \koff u_b(r) \right] u_{b, -}(r) \d r \\
        &= 2 \int_{0}^{t} \left[ - \delta_{u_b} \left( u_{b,-}(r) \right)^2 - \kon u^2(r) \gamma v^*(r, x) u_{b,-}(r) - \koff \left( u_{b,-}(r) \right)^2 \right] \d r \leq 0,
    \end{align*}

    \noindent
    so $u_{b,-}(t) = 0$ for every $t \in [0, \beta]$, showing that $u_b$ must be non-negative for almost every $t \in [0, \beta]$. Using this in the estimate for $u$, we get
%
    \begin{align*}
        \left( u_{-}(t) \right)^2 &= 2 \int_{0}^{t} u_{-}(r) \frac{\d}{\d t} u_{-}(r)\d r = - 2 \int_{0}^{t} u_{-}(r) u_{t}(r) \d r \\
        &= - 2 \int_{0}^{t} \left[ P_{u} + \alpha u_b(r) - \delta_{u} u(r) - 2 \kon u(r)^2 \gamma v^*(r, x) + 2 \koff u_b(r) \right] u_{-}(r) \d r \\
        &= 4 \kon \int_{0}^{t} u(r)^2 \gamma v^*(r, x)u_{-}(r) \d r \\
        & \qquad - 2 \int_{0}^{t} \left[ P_{u} u_{-}(r) + \alpha u_b(r) u_{-}(r) + \delta_{u} u(r)^2 + 2 \koff u_b(r)u_{-}(r) \right] \d r \\
        &\leq 4 \kon \int_{0}^{t} u(r)^2 \gamma v^*(r, x)u_{-}(r) \d r.
    \end{align*}

    \noindent
    By application of Gronwall's lemma, this shows that $u_-(t) = 0$ for almost all $t \in [0, \beta]$, so $u \geq 0$.
    
    We can use the non-negativity of $u$ and $u_b$ to show that these functions are bounded. Define $r := u + 2 u_b$ and note that
%
    \begin{equation*}
        r_t = u_t + 2 u_{b, t} = P_u - \delta_u u + (\alpha - 2 \delta_{u_b}) u_b \leq P_u + m r,
    \end{equation*}

    \noindent
    where $m := \max\left(\delta_u, \frac{1}{2}(\alpha - 2 \delta_{u_b})\right)$. Then Gronwall's lemma gives
%
    \begin{equation*}
        r(t) \leq e^{mt}\left(r(0) + t P_u\right).
    \end{equation*}

    \noindent
    Now the non-negativity of both $u$ and $u_b$ and the essential boundedness of $u_0$ and $u_{b,0}$ allows us to bound the solutions independent of $\widetilde{T}$ by
%
    \begin{equation}\label{eq:ode_bounds}
    \begin{alignedat}{2}
        u &\leq e^{mT} \left(M + P_u\right) &&=: M_u(T) \qquad \text{and} \\
        u_b &\leq \frac{1}{2} M_u(T) &&=: M_{u_b}(T),
    \end{alignedat}
    \end{equation}

    \noindent
    where $M = \esssup(u_0 + 2 u_{b, 0})$, for almost every $x \in \Gamma$ and $t \in [0, \beta]$.

    Now suppose we have solutions $u$ and $u_b$ on a maximum time interval $[0, \widetilde{\beta})$. Since we have a bound for $u$ and $u_b$ that is independent of $\widetilde{\beta}$, we can use the continuation argument from \cite{coddington_theory_1984} and see that we have a solution for $u$ and $u_b$ on $[0, \widetilde{T}]$.

    Now if we let $x$ vary, we can see that $u$ and $u_b$ are both of $L^2$ regularity in space, because the bounds from \cref{eq:ode_bounds} are independent of space. Lastly, the fact that $u$ and $u_b$ are essentially bounded shows that $u_t$ and $u_{b, t}$ are bounded, so $u, u_b \in H^1(0, \widetilde{T}; \Ltwo{\Gamma})$.
\end{proof}

\endgroup

We can use a similar approach for the PDE in \cref{sys:PDEs}. Fix $(u^*, u_b^*) \in \odespace$ and consider the modified weak formulation given by
%
\begin{equation}\label{eq:frozen_pde}
\begin{split}
    \pair{v_t, \phi} = \inner{P_v, \phi}_{\Ltwo{\Gamma}} - \kon \inner{\left(u^*\right)^2 v, \phi}_{\Ltwo{\Gamma}} + \koff \inner{u_b^*, \phi}_{\Ltwo{\Gamma}} \\
    - D_v \inner{\nabla v, \nabla \phi}_{\Ltwo{\Omega; \R^d}} - \delta_v \inner{v, \phi}_{\Ltwo{\Omega}}.
\end{split}
\end{equation}

\noindent
for almost every $t \in [0, \widetilde{T}]$. This PDE is now decoupled from the ODEs, and we can prove the existence and uniqueness independently from the ODE system.

\begin{lemma}{Existence and uniqueness for the decoupled PDE}{exun_frozen_pde}
    Let the initial condition $v_0 \in \Ltwo{\Omega}$ be essentially bounded and non-negative and fix $(u^*, u_b^*) \in \odespace$. Then there exists a unique weak solution $v \in \pdespace$ to \cref{eq:frozen_pde} satisfying these initial conditions.
\end{lemma}

\begin{proof}
    To prove the existence and uniqueness goes beyond the scope of this thesis. Hence, we assume that a unique weak solution $v$ of the PDE exists and $v$ is essentially bounded by a constant depending on $T$. This can be proven using Galerkin methods, as demonstrated in \cite{evans2010partial}. Let $v$ now be such a solution. To show that $v$ is non-negative, we estimate the negative part of $v$. Let $v_{+}$ and $v_{-}$ be defined as $v_+ := \max(0, v)$ and $v_- := - \min(0, v)$ respectively, such that $v = v_{+} - v_{-}$. Then as discussed in \cite{evans2010partial}, we have $v_{+}, v_{-} \in \Ltwo{0, \widetilde{T}; H^1(\Omega)}$ and $\left(v_{+}\right)_t, \left(v_{-}\right)_t \in \Ltwo{0, \widetilde{T}; H^1(\Omega)^*}$. Furthermore, $v_t v_{-} = - (v_{-})_t v_{-}$, as well as $\nabla v \cdot \nabla v_{-} = - \left|\nabla v_{-}\right|^2$, so for all almost every $t$ we have
%
    \begin{equation*}
         \frac{\d}{\d t} \frac12 \norm{v_{-}}_{\Ltwo{\Omega}}^2 = \pair{(v_-)_t, v_-} = - \pair{v_t, v_{-}}
    \end{equation*}

    \noindent
    due to \cref{thm:lionsmagenes}. We use $v_-$ as a test function and employ the fact that $v$ is a solution to the PDE, which gives
    %
    \begin{equation*}
        \begin{aligned}
         \frac{\d}{\d t} &\frac12 \norm{v_{-}}_{\Ltwo{\Omega}}^2 = - \pair{v_t, v_{-}} = D_v \int_{\Omega} \nabla v \cdot \nabla v_- + \delta_v \int_{\Omega} v v_- - \int_{\Gamma} \left(P_v - \kon u^2 v + \koff u_b\right) v_-  \d S \\
        &= - D_v \norm{\nabla v_-}_{\Ltwo{\Omega; \R^d}}^2 - \delta_v \norm{v_-}_{\Ltwo{\Omega}}^2 - P_v \int_{\Gamma} v_- \d S + \kon \int_{\Gamma} u^2 v v_-  \d S - \koff \int_{\Gamma} u_b v_-  \d S \\
        &= - D_v \norm{\nabla v_-}_{\Ltwo{\Omega; \R^d}}^2 - \delta_v \norm{v_-}_{\Ltwo{\Omega}}^2 - P_v \int_{\Gamma} v_- \d S - \kon \int_{\Gamma} u^2 (v_-)^2  \d S - \koff \int_{\Gamma} u_b v_-  \d S \\
        &\leq 0,
    \end{aligned}
    \end{equation*}

    \noindent
    where we used the non-negativity of $v_-$ and $u_b$ in the last step. By \cref{thm:lionsmagenes}, we know that $v$ is continuous in time with respect to the $L^2$-norm. The only way for this to be achieved while also guaranteeing the time derivative is non-positive, is by having
%
    \begin{equation*}
        \norm{v_{-}}_{\Ltwo{\Omega}} = 0
    \end{equation*}

    \noindent
    for almost every $t$, which implies that $v$ is non-negative.

    % To show boundedness of $v$, we consider $\left( v - M_1 e^{M_2 t} \right)_+$ as a test function for constants $M_1$ and $M_2$. This gives

    % \begin{align*}
    %     \frac{1}{2} \frac{\d}{\d t} \norm{\left( v - M_1 e^{M_2 t} \right)_+}&_{\Ltwo{\Omega}}^2 = \pair{\frac{\d}{\d t} \left( v - M_1 e^{M_2 t} \right)_+, \left( v - M_1 e^{M_2 t} \right)_+} \\
    %     =& \pair{\left( v_t - M_1 M_2 e^{M_2 t} \right), \left( v - M_1 e^{M_2 t} \right)_+} \\
    %     =& \int_\Gamma \left(P_v - \kon \left(u^*\right)^2 v + \koff \left(u_b^*\right)\right) \left( v - M_1 e^{M_2 t} \right)_+ \d S \\
    %      &- \delta_v \int_{\Omega} \left( v - M_1 M_2 e^{M_2 t} \right) \left( v - M_1 e^{M_2 t} \right)_+ \\
    % \end{align*}

    % \todo{Mariya's term?}

    % {\color{red}To show that the map $t \mapsto \norm{v(t)}_{\Ltwo{\Omega}}^2$ is essentially bounded, we use $v$ itself as a test function to get

    % \begin{equation*}
    % \begin{aligned}
    %      \frac{\d}{\d t} \frac12 &\norm{v}_{\Ltwo{\Omega}}^2 = \int_{\Omega} v_t v = \int_{\Gamma} \left(P_v - \kon u^2 v + \koff u_b\right) v \d S - D_v \int_{\Omega} \nabla v \cdot \nabla v - \delta_v \int_{\Omega} v v \\
    %     &= P_v \int_{\Gamma} v \d S - \kon \int_{\Gamma} u^2 v^2 \d S + \koff \int_{\Gamma} u_b v \d S  - D_v \norm{\nabla v}_{\Ltwo{\Omega; \R^d}}^2 - \delta_v \norm{v}_{\Ltwo{\Omega}}^2 \\
    %     &\leq P_v \int_{\Gamma} v \d S + \koff \int_{\Gamma} u_b v \d S - D_v \norm{\nabla v}_{\Ltwo{\Omega; \R^d}}^2 . \\
    % \end{aligned}
    % \end{equation*}

    % \noindent
    % Now we apply Young's equality with $\eps$ to the first two integrals in this last expression to derive   

    % \begin{equation*}
    % \begin{alignedat}{2}
    %     P_v \int_{\Gamma} v \d S &\leq \int_{\Gamma} \left(\eps v^2 + \frac{1}{4 \eps} P_v^2 \right) \d S = \eps \norm{v}_{\Ltwo{\Gamma}}^2 + \frac{1}{4\eps} P_v^2 | \Gamma | &&=: \eps \norm{v}_{\Ltwo{\Omega}}^2 + \frac{1}{2}B_\eps ,\\
    %     \koff \int_\Gamma u_b v &\leq \int_\Gamma \left(\eps + \frac{1}{4 \eps} \koff^2 u_b^2 \right) \d S = \eps \norm{v}_{\Ltwo{\Gamma}}^2 + \frac{1}{4\eps} \koff^2 \norm{w}_{\Ltwo{\Gamma}}^2 &&=: \eps \norm{v}_{\Ltwo{\Gamma}}^2 + \frac{1}{2} D_\eps \norm{w}_{\Ltwo{\Gamma}}^2 ,
    % \end{alignedat}
    % \end{equation*}

    % \noindent
    % which, when plugged into the original estimate, gives

    % \begin{equation*}
    %      \frac{\d}{\d t} \frac12 \norm{v}_{\Ltwo{\Omega}}^2 \leq \frac{1}{2} B_\eps + \frac{1}{2} D_\eps \norm{w}_{\Ltwo{\Gamma}}^2 + 2 \eps \norm{v}_{\Ltwo{\Gamma}}^2 - D_v \norm{\nabla v}_{\Ltwo{\Omega; \R^d}}^2 .
    % \end{equation*}

    % \noindent
    % To convert the $\Ltwo{\Gamma}$-norm of $v$ on the right-hand side into an $\Ltwo{\Omega}$-norm, we employ the trace inequality from \cref{thm:trace} and get

    % \begin{equation*}
    % \begin{aligned}
    %      \frac{\d}{\d t} \frac12 \norm{v}_{\Ltwo{\Omega}}^2 &\leq \frac{1}{2} B_\eps + \frac{1}{2} D_\eps \norm{w}_{\Ltwo{\Gamma}}^2 + 2 \eps C_\gamma^2 \norm{v}_{H^1(\Omega)}^2 - D_v \norm{\nabla v}_{\Ltwo{\Omega; \R^d}}^2 \\
    %     &= \frac{1}{2} B_\eps + \frac{1}{2} D_\eps \norm{w}_{\Ltwo{\Gamma}}^2 + 2 \eps C_\gamma^2 \norm{v}_{\Ltwo{\Omega}}^2 + 2 \eps C_\gamma^2 \norm{\nabla v}_{\Ltwo{\Omega; \R^d}}^2 - D_v \norm{\nabla v}_{\Ltwo{\Omega; \R^d}}^2 \\
    %     &\leq \frac{1}{2} B_\eps + \frac{1}{2} D_\eps \norm{w}_{\Ltwo{\Gamma}}^2 + 2 \eps C_\gamma^2 \norm{v}_{\Ltwo{\Omega}}^2 \\
    %     &=: \frac{1}{2} B_\eps + \frac{1}{2} D_\eps \norm{w}_{\Ltwo{\Gamma}}^2 + \frac{1}{2} A_\eps \norm{v}_{\Ltwo{\Omega}}^2
    % \end{aligned}
    % \end{equation*}

    % \noindent
    % for small enough $\eps > 0$ and almost every $t$. We are now ready to use Gronwall's lemma to get the final estimate \todo{gronwall application}

    % \begin{equation*}
    % \begin{aligned}
    %     \norm{v(t)}_{\Ltwo{\Omega}}^2 &\leq \norm{v_0}_{\Ltwo{\Omega}}^2 e^{A_\eps t} + \int_{0}^{t} \left(B_\eps + D_\eps \norm{w(s)}_{\Ltwo{\Gamma}}^2\right) e^{A_\eps (t - s)} \d s \\
    %     &= \norm{v_0}_{\Ltwo{\Omega}}^2 e^{A_\eps t} + \frac{B_\eps}{A_\eps} \left(e^{A_\eps t} - 1\right) + D_\eps \int_{0}^{t} \norm{w(s)}_{\Ltwo{\Gamma}}^2 e^{A_\eps (t - s)} \d s \\
    %     &\leq \norm{v_0}_{\Ltwo{\Omega}}^2 e^{A_\eps t} + \frac{B_\eps}{A_\eps} \left(e^{A_\eps t} - 1\right) + \frac{D_\eps M_{u_b}(T)^2 |\Gamma|}{A_\eps} \left(e^{A_\eps t} - 1\right) \\
    %     &\leq \norm{v_0}_{\Ltwo{\Omega}}^2 e^{A_\eps T} + \frac{B_\eps}{A_\eps} \left(e^{A_\eps T} - 1\right) + \frac{D_\eps M_{u_b}(T)^2 |\Gamma|}{A_\eps} \left(e^{A_\eps T} - 1\right) \\
    %     &=: M_v(T)
    % \end{aligned}
    % \end{equation*}

    % \noindent
    % for almost every $t$, where $M_{u_b}(T)$ was defined in the proof of \cref{lem:exun_frozen_odes}.}
\end{proof}

Using \cref{lem:exun_frozen_odes,lem:exun_frozen_pde}, we can define two operators between $\pdespace$ and $\odespace$. One takes a function $v^* \in \pdespace$ and maps it to the unique solution $(u, u_b) \in \odespace$ of \cref{eq:frozen_odes}. Let us refer to this operator as $\mathcal{T}_{1}: \pdespace \to \odespace$. The other operator, $\mathcal{T}_{2}: \odespace \to \pdespace$, takes two arbitrary functions $(u^*, u_b^*) \in \odespace$ and defines $\mathcal{T}_{2}(u^*, u_b^*)$ as the unique solution $v \in \pdespace$ to \cref{eq:frozen_pde}. Combining these two operators, we define  
%
\begin{equation*}
    \mathcal{T} := \mathcal{T}_{2} \circ \mathcal{T}_{1} : \pdespace \to \pdespace.
\end{equation*}

\noindent
We want to show that this operator has a unique fixed point. If this is the case, we have a functions $v \in \pdespace$ such that $\mathcal{T}(v) = v$, and let $(u, u_b) \in \odespace$ denote $(u, u_b) = \mathcal{T}_{1}(v)$. Given that then $\mathcal{T}_{2}(u, u_b) = v$, by definition this means that $v$ solves \cref{eq:weak_pde} for almost every $t \in (0, \widetilde{T})$, and $w = (u, u_b)$ solves \cref{eq:weak_ode} for almost every $t \in (0, \widetilde{T})$ and $x \in \Gamma$, meaning that $v$, $u$ and $u_b$ are weak solutions of our model on the local time scale $(0, \widetilde{T})$. Conversely, if $v$, $u$ and $u_b$ are weak solutions to the model, then by the uniqueness of solutions in \cref{lem:exun_frozen_odes,lem:exun_frozen_pde}, we must have that $\mathcal{T}_{2}(u, u_b) = v$ and $\mathcal{T}_{1}(v) = (u, u_b)$, and thus $\mathcal{T}(v) = v$.

It should be clear now how we will apply Banach's fixed-point theorem. If we prove that $\mathcal{T}$ is a contraction, we can guarantee the existence of a fixed point and therefore a weak solution. Since every weak solution of the model is also a fixed point, the theorem also grants us uniqueness of such solutions. To prove that $\mathcal{T}$ is a contraction, we first show this for $\mathcal{T}_{2}$ and $\mathcal{T}_{1}$ separately, after which we can easily combine the results.

\begingroup
\newcommand{\vb}{\overline{v}}
\newcommand{\vt}{\widetilde{v}}
\newcommand{\ub}{\overline{u}}
\newcommand{\ut}{\widetilde{u}}
\newcommand{\ubb}{\overline{u}_b}
\newcommand{\ubt}{\widetilde{u}_b}

\begin{lemma}{}{contraction_pde}
    For $\widetilde{T} > 0$ small enough, $\mathcal{T}_{2}: \odespace \to \pdespace$ is a contraction. That is, there is a $q \in [0, 1)$ such that
    %
    \begin{equation*}
        \norm{\mathcal{T}_{2} (\ub, \ubb) - \mathcal{T}_{2} (\ut, \ubt)}_{\Ltwo{0, \widetilde{T}; \Ltwo{\Omega}}}^2 \leq q \norm{(\ub, \ubb) - (\ut, \ubt)}_{\left(\Ltwo{0, \widetilde{T}; \Ltwo{\Gamma}}\right)^2}^2
    \end{equation*}

    for every $(\ub, \ubb), (\ut, \ubt) \in \odespace$.
\end{lemma}

\begin{proof}
    Let $(\ub, \ubb), (\ut, \ubt) \in \odespace$ and define $\vb, \vt \in \pdespace$ by $\vb := \mathcal{T}_{2}(\ub, \ubb)$ and $\vt := \mathcal{T}_{2}(\ut, \ubt)$. We now denote their differences by
%
    \begin{equation*}
        v := \vb - \vt, \qquad u := \ub - \ut \qquad \text{and} \qquad u_b := \ubb - \ubt,
    \end{equation*}

    \noindent
    such that
%
    \begin{equation*} 
        \norm{\mathcal{T}_{2} (\ub, \ubb) - \mathcal{T}_{2} (\ut, \ubt)}_{\Ltwo{0, \widetilde{T}; \Ltwo{\Omega}}}^2 = \norm{\vb - \vt}_{\Ltwo{0, \widetilde{T}; \Ltwo{\Omega}}}^2 = \norm{v}_{\Ltwo{0, \widetilde{T}; \Ltwo{\Omega}}}^2.
    \end{equation*}

    \noindent
    We can estimate this norm using the weak form of the PDE. We use \cref{thm:lionsmagenes} to see that
%
    \begin{equation*}
    \begin{aligned}
        \frac{\d}{\d t} \frac{1}{2} \norm{v}_{\Ltwo{\Omega}}^2 = \pair{v_t, v} = - D_v \int_\Omega \left|\nabla v\right|^2 - \delta_v \int_\Omega v^2 + \int_\Gamma \left(- \kon \ub^2 \vb + \koff \ubb + \kon \ut^2 \vt - \koff \ubt\right) v \d S ,
    \end{aligned}
    \end{equation*}

    \noindent
    which after rearranging becomes
%
    \begin{equation*}
    \begin{aligned}
        \frac{\d}{\d t} \frac{1}{2} \norm{v}_{\Ltwo{\Omega}}^2 + D_v \norm{\nabla v}_{\Ltwo{\Omega; \R^d}}^2 + \delta_v \norm{v}_{\Ltwo{\Omega}}^2 &= - \kon \int_\Gamma \left(\ub^2 \vb - \ut^2 \vt\right)v \d S + \koff \int_\Gamma u_b v \d S
    \end{aligned}
    \end{equation*}

    \noindent
    for almost all $t$. We then use the identity
%
    \begin{equation} \label{eq:contraction_identity_1}
        \ub^2 \vb - \ut^2 \vt = \ub^2 \vb - \ut^2 \vb + \ut^2 \vb - \ut^2 \vt = (\ub^2 - \ut^2) \vb + \ut^2 v = (\ub + \ut) \vb u + \ut^2 v
    \end{equation}

    \noindent
    to estimate
%
    \begin{equation*}
    \begin{aligned}
        \frac{\d}{\d t} \frac{1}{2} \norm{v}_{\Ltwo{\Omega}}^2 + &D_v \norm{\nabla v}_{\Ltwo{\Omega; \R^d}}^2 + \delta_v \norm{v}_{\Ltwo{\Omega}}^2 \\
        &= - \kon \int_\Gamma \left(\ub + \ut\right) \vb u v \d S - \kon \int_\Gamma \ut^2 v^2 \d S + \koff \int_\Gamma u_b v \d S \\
        &\leq \kon \int_\Gamma \left|(\ub + \ut) \vb u v\right| \d S + \koff \int_\Gamma u_b v \d S \\
        &\leq 2 \kon M_{u}(T) M_v(T) \int_\Gamma \left|u v\right| \d S + \koff \int_\Gamma u_b v \d S .\\
    \end{aligned}
    \end{equation*}

    \noindent
    The mixed terms can be separated using Young's inequality, giving
%
    \begingroup
    \allowdisplaybreaks
    \begin{align*}
        \frac{\d}{\d t} \frac{1}{2} &\norm{v}_{\Ltwo{\Omega}}^2 + D_v \norm{\nabla v}_{\Ltwo{\Omega; \R^d}}^2 + \delta_v \norm{v}_{\Ltwo{\Omega}}^2 \\
        &\leq 2 \kon M_{u}(T) M_v(T) \left(A_{\eps} \norm{u}_{\Ltwo{\Gamma}}^2 + \eps \norm{v}_{\Ltwo{\Gamma}}^2\right) + \koff \left(A_{\eps} \norm{u_b}_{\Ltwo{\Gamma}}^2 + \eps \norm{v}_{\Ltwo{\Gamma}}^2\right) \\
        &\leq B_\eps(T) \norm{u}_{\Ltwo{\Gamma}}^2 + \koff A_{\eps} \norm{u_b}_{\Ltwo{\Gamma}}^2 + C(T) \eps \norm{v}_{\Ltwo{\Gamma}}^2 \\
        &\leq B_\eps(T) \norm{u}_{\Ltwo{\Gamma}}^2 + \koff A_{\eps} \norm{u_b}_{\Ltwo{\Gamma}}^2 + C_\gamma^2 C(T) \eps \norm{v}_{H^1(\Omega)}^2\\
    \end{align*}
    \endgroup

    \noindent
    for almost all $t$. Now let $\delta := \min (\delta_v, D_v)$ and integrate from 0 to $\widetilde{T}$ to get
%
    \begin{equation*}
        \frac{1}{2} \norm{v(t)}_{\Ltwo{\Omega}}^2 + (\delta - C_\gamma^2 C(T) \eps) \norm{v}_{\Ltwo{0, \widetilde{T}; H^1(\Omega)}}^2 \leq B_\eps(T) \norm{u}_{\Ltwo{0, \widetilde{T}; \Ltwo{\Gamma}}}^2 + \koff A_{\eps} \norm{u_b}_{\Ltwo{0, \widetilde{T}; \Ltwo{\Gamma}}}^2
    \end{equation*}

    \noindent
    for almost all $t$. If we now choose $\eps < \frac{\delta}{C_\gamma^2 C(T)}$, we can estimate both the $\Ltwo{0, \widetilde{T}; H^1(\Omega)}$-norm and the $\Ltwo{0, \widetilde{T}; \Ltwo{\Omega}}$-norm of $v$. The first will only be helpful later, but the latter concludes this proof. First, since $\delta - C_\gamma^2 C(T) \eps > 0$, we have
%
    \begin{equation}\label{eq:extra_estimate}
    \begin{aligned}
        \norm{v}_{\Ltwo{0, \widetilde{T}; H^1(\Omega)}}^2 \leq D_\eps(T) \left(\norm{u}_{\Ltwo{0, \widetilde{T}; \Ltwo{\Gamma}}}^2 + \norm{u_b}_{\Ltwo{0, \widetilde{T}; \Ltwo{\Gamma}}}^2\right) .
    \end{aligned}
    \end{equation}

    \noindent
    Second, we can integrate from 0 to $\widetilde{T}$ once more to get
%
    \begin{equation*}
        \norm{v}_{\Ltwo{0, \widetilde{T}; \Ltwo{\Omega}}}^2 \leq \widetilde{T} M_\eps(T) \left(\norm{u}_{\Ltwo{0, \widetilde{T}; \Ltwo{\Gamma}}}^2 + \norm{u_b}_{\Ltwo{0, \widetilde{T}; \Ltwo{\Gamma}}}^2\right).
    \end{equation*}

    \noindent
    Now for $\widetilde{T} < \frac{1}{M_\eps(T)}$, this coefficient is less than 1, so $\mathcal{T}_2$ will be a contraction.
\end{proof}

\begin{lemma}{}{contraction_odes}
    For $\widetilde{T} > 0$ small enough, $\mathcal{T}_{1}: \pdespace \to \odespace$ is a contraction. That is, there is a $q \in [0, 1)$ such that
    %
    \begin{equation*}
        \norm{\mathcal{T}_{1} (\vb) - \mathcal{T}_{1} (\vt)}_{\left(\Ltwo{0, \widetilde{T}; \Ltwo{\Gamma}}\right)^2}^2 \leq q \norm{\vb - \vt}_{\Ltwo{0, \widetilde{T}; \Ltwo{\Omega}}}^2
    \end{equation*}

    for every $\vb, \vt \in \pdespace$.
\end{lemma}

\begin{proof}
    Now let $\vb, \vt \in \pdespace$ and define $\ub, \ubb, \ut$ and $\ubt$ by $(\ub, \ubb) := \mathcal{T}_{1}(\vb)$ and $(\ut, \ubt) := \mathcal{T}_{1}(\vt)$. Once again, we denote their differences by
%
    \begin{equation*}
        v := \vb - \vt, \qquad u := \ub - \ut \qquad \text{and} \qquad u_b := \ubb - \ubt,
    \end{equation*}

    \noindent
    such that
%
    \begin{align*}
        \norm{\mathcal{T}_{1} (\vb) - \mathcal{T}_{1} (\vt)}_{\left(\Ltwo{0, \widetilde{T}; \Ltwo{\Gamma}}\right)^2}^2 &= \norm{(\ub, \ubb) - (\ut, \ubt)}_{\left(\Ltwo{0, \widetilde{T}; \Ltwo{\Gamma}}\right)^2}^2 = \norm{(u, u_b)}_{\left(\Ltwo{0, \widetilde{T}; \Ltwo{\Gamma}}\right)^2}^2 \\
        &= \norm{u}_{\Ltwo{0, \widetilde{T}; \Ltwo{\Gamma}}}^2 + \norm{u_b}_{\Ltwo{0, \widetilde{T}; \Ltwo{\Gamma}}}^2.
    \end{align*}

    \noindent
    We therefore seek to estimate these separate differences $u$ and $u_b$ in terms of $v$. Starting off with $u$, we see that it satisfies
    \begin{equation*}
    \begin{aligned}
        \frac{\d}{\d t} \norm{u}_{\Ltwo{\Gamma}}^2 &= 2 \int_\Gamma u_t u \d S = 2 \int_\Gamma (\ub_t - \ut_t) u \d S \\
        &= 2 \int_\Gamma \left( \alpha (\ubb - \ubt) - \delta_u (\ub - \ut) - 2 \kon (\ub^2 \vb - \ut^2 \vt) + 2 \koff (\ubb - \ubt) \right) u \d S \\
        &= 2 (\alpha + 2 \koff) \int_\Gamma u u_b \d S - 2 \delta_u \int_\Gamma u^2 \d S - 4 \kon \int_\Gamma (\ub^2 \vb - \ut^2 \vt) u \d S
    \end{aligned}
    \end{equation*}

    \noindent
    for almost every $t$. Here, the first step is justified because we know that $u$ belongs to $H^1(0, \widetilde{T}; \Ltwo{\Gamma})$. As we did in the proof of \cref{lem:contraction_pde}, we now use the identity in \cref{eq:contraction_identity_1} to conclude that
    %
    \begin{equation*}
    \begin{aligned}
        \frac{\d}{\d t} \norm{u}_{\Ltwo{\Gamma}}^2 &= 2 (\alpha + 2 \koff) \int_\Gamma u u_b \d S - 2 \delta_u \int_\Gamma u^2 \d S - \overbrace{4 \kon \int_\Gamma (\ub + \ut) \vb u^2 \d S }^{\geq 0} - 4 \kon \int_\Gamma \ut^2 v u \d S \\
        &\leq 2 (\alpha + 2 \koff) \int_\Gamma u u_b \d S - 2 \delta_u \int_\Gamma u^2 \d S + 4 \kon \int_\Gamma \ut^2 \left| v u \right| \d S \\
        &\leq 2(\alpha + 2 \koff) \int_\Gamma u u_b \d S - 2 \delta_u \norm{u}_{\Ltwo{\Gamma}}^2 + 4 \kon M_u (T)^2 \int_\Gamma | v u | \d S . \\
    \end{aligned}
    \end{equation*}
    
    \noindent
    To further separate the mixed terms, we use Young's inequality twice and obtain the estimate
%
    \begin{equation*}
    \begin{aligned}
        \frac{\d}{\d t} \norm{u}_{\Ltwo{\Gamma}}^2 &\leq A_1(T) \norm{u}_{\Ltwo{\Gamma}}^2 + A_2 \norm{u_b}_{\Ltwo{\Gamma}}^2 + \widetilde{B_1}(T) \norm{v}_{\Ltwo{\Gamma}}^2 , \\
        &\leq A_1(T) \norm{u}_{\Ltwo{\Gamma}}^2 + A_2 \norm{u_b}_{\Ltwo{\Gamma}}^2 + B_1(T) \norm{v}_{H^1(\Omega)}^2 , \\
    \end{aligned}
    \end{equation*}

    \noindent
    where $A_1(T)$, $A_2$, $\widetilde{B_1}(T)$ and $B_1(T)$ are defined as
%
    \begin{equation*}
        \begin{aligned}
            A_1(T) &:= - 2 \delta_u + \alpha + 2 \koff + 2 \kon M_u (T)^2 ,& A_2 &:= \alpha + 2 \koff \\
            \widetilde{B_1}(T) &:= 2 \kon M_u (T)^2 \qquad \text{and} & B_1(T) &:= C_\gamma^2 \widetilde{B_1}(T) .
        \end{aligned}
    \end{equation*}

    \noindent
    Using the same arguments, we get a similar estimate for $\frac{\d}{\d t} \norm{u_b}_{\Ltwo{\Gamma}}^2$, namely

    \begingroup
    \allowdisplaybreaks
    \begin{align*}
        \frac{\d}{\d t} \norm{u_b}_{\Ltwo{\Gamma}}^2 &= 2 \int_\Gamma \left(\left(\ubb\right)_t - \left(\ubt\right)_t\right) u_b \d S \\
        &= - 2 \left(\delta_{u_b} + \koff\right) \int_\Gamma u_b^2 \d S + 2 \kon \int_\Gamma \left(\ub^2 \vb - \ut^2 \vt\right) u_b \d S \\
        &= - 2 \left(\delta_{u_b} + \koff\right) \int_\Gamma u_b^2 \d S + 2 \kon \int_\Gamma \left(\ub + \ut\right) \vt u u_b \d S + 2 \kon \int_\Gamma \ut^2 v u_b \d S \\
        &\leq - 2 \left(\delta_{u_b} + \koff\right) \int_\Gamma u_b^2 \d S + 2 \kon \int_\Gamma \left(\ub + \ut\right) \vt \left| u u_b \right| \d S + 2 \kon \int_\Gamma \ut^2 \left| v u_b \right| \d S \\
        &\leq - 2 \left(\delta_{u_b} + \koff\right) \int_\Gamma u_b^2 \d S + 4 \kon M_u(T) M_v(T) \int_\Gamma \left| u u_b \right| \d S + 2 \kon M_u(T)^2 \int_\Gamma \left| v u_b \right| \d S \\
        &\leq A_3(T) \norm{u_b}_{\Ltwo{\Gamma}}^2 + A_4(T) \norm{u}_{\Ltwo{\Gamma}}^2 + \widetilde{B_2}(T) \norm{v}_{\Ltwo{\Gamma}}^2 ,\\
        &\leq A_3(T) \norm{u_b}_{\Ltwo{\Gamma}}^2 + A_4(T) \norm{u}_{\Ltwo{\Gamma}}^2 + B_2(T) \norm{v}_{H^1(\Omega)}^2 ,\\
    \end{align*}
    \endgroup

    \noindent
    where $A_3(T)$, $A_4(T)$, $\widetilde{B_2}(T)$ and $B_2(T)$ are given by
%
    \begin{equation*}
        A_3(T) := 2 \kon M_u(T) M_v(T) + \kon M_u(T)^2 - 2 \delta_{u_b} - 2 \koff ,
    \end{equation*}
    \begin{equation*}
        A_4(T) := 2 \kon M_u(T) M_v(T), \qquad \widetilde{B_2}(T) := \kon M_u(T)^2  \qquad \text{and} \qquad B_2(T) := C_\gamma^2 \widetilde{B_2}(T) .
    \end{equation*}
    
    \noindent
    To combine these estimates, we set
    %
    \begin{equation*}
        A(T) := 2 \max \left(A_1(T), A_2, A_3(T), A_4(T)\right) \qquad \text{and} \qquad B(T) := 2 \max \left(B_1(T), B_2(T)\right)
    \end{equation*}

    \noindent
    and sum them to get
%
    \begin{equation*}
        \frac{\d}{\d t} \left(\norm{u}_{\Ltwo{\Gamma}}^2 + \norm{u_b}_{\Ltwo{\Gamma}}^2\right) \leq A(T) \left(\norm{u}_{\Ltwo{\Gamma}}^2 + \norm{u_b}_{\Ltwo{\Gamma}}^2\right) + B(T) \norm{v}_{H^1(\Omega)}^2
    \end{equation*}

    \noindent
    for almost all $t$. An application of Gronwall's lemma gives
    %
    \begin{equation*}
    \begin{aligned}
        \norm{u(t)}_{\Ltwo{\Gamma}}^2 + \norm{u_b(t)}_{\Ltwo{\Gamma}}^2 &\leq e^{t A(T)} \left[ \norm{u(0)}_{\Ltwo{\Gamma}}^2 + \norm{u_b(0)}_{\Ltwo{\Gamma}}^2 + B(T) \int_{0}^{t} \norm{v(s)}_{H^1(\Omega)}^2 \d s \right] \\
        &\leq B(T) e^{T A(T)} \norm{v}_{\Ltwo{0, \widetilde{T}; H^1(\Omega)}}^2 , \\
    \end{aligned}
    \end{equation*}

    \noindent
    since the initial conditions of $\ub$ and $\ut$ are the same, and therefore $u(0) = 0$ for every $x \in \Gamma$. Naturally, the same holds for $u_b(0)$. After integrating from 0 to $\widetilde{T}$, we can conclude that
%
    \begin{equation*}
        \norm{u}_{\Ltwo{0, \widetilde{T}; \Ltwo{\Gamma}}}^2 + \norm{u_b}_{\Ltwo{0, \widetilde{T}; \Ltwo{\Gamma}}}^2 \leq \widetilde{T} B(T) e^{T A(T)} \norm{v}_{\Ltwo{0, \widetilde{T}; H^1(\Omega)}}^2 .
    \end{equation*}

    \noindent
    It may seem like we are done now, but the current estimate uses the wrong norm on $v$. To get an estimate using the $\Ltwo{0, \widetilde{T}; \Ltwo{\Omega}}$-norm instead, we can use \cref{eq:extra_estimate} to get

    \begin{multline}
        \norm{u}_{\Ltwo{0, \widetilde{T}; \Ltwo{\Gamma}}}^2 + \norm{u_b}_{\Ltwo{0, \widetilde{T}; \Ltwo{\Gamma}}}^2 \\
        \leq \widetilde{T} B(T) e^{T A(T)} D_\eps(T) \left(\norm{u}_{\Ltwo{0, \widetilde{T}; \Ltwo{\Gamma}}}^2 + \norm{u_b}_{\Ltwo{0, \widetilde{T}; \Ltwo{\Gamma}}}^2\right) + \overbrace{\widetilde{T} B(T) e^{T A(T)} \norm{v}_{\Ltwo{0, \widetilde{T}; \Ltwo{\Omega}}}^2}^{\geq 0},
    \end{multline}

    \noindent
    or after rearranging terms,
%
    \begin{equation*}
    \begin{aligned}
        \norm{u}_{\Ltwo{0, \widetilde{T}; \Ltwo{\Gamma}}}^2 + \norm{u_b}_{\Ltwo{0, \widetilde{T}; \Ltwo{\Gamma}}}^2 \leq \frac{\widetilde{T} B(T) e^{T A(T)}}{1 - \widetilde{T} B(T) e^{T A(T)} D_\eps(T)} \norm{v}_{\Ltwo{0, \widetilde{T}; \Ltwo{\Omega}}}^2.
    \end{aligned}
    \end{equation*}

    \noindent
    As this factor clearly approaches 0 when $\widetilde{T}$ goes to 0, we can choose $\widetilde{T}$ sufficiently small such that it is less than 1 and conclude the proof.
\end{proof}

We have now gathered the needed auxiliary results to prove the main statement of this section.

\begin{proof}[Proof of \cref{thm:eq_and_un}]
    Let $\widetilde{T} > 0$ be small enough such that, by \cref{lem:contraction_pde,lem:contraction_odes}, $\mathcal{T}_1$ and $\mathcal{T}_2$ are contractions with Lipschitz constants $q_1$ and $q_2$, respectively. Define, as discussed before, $\mathcal{T}$ as the composition of $\mathcal{T}_1$ and $\mathcal{T}_2$. Then for any $\vb, \vt \in \Ltwo{0, \widetilde{T}; \Ltwo{\Omega}}$, we have
%
    \begin{equation*}
        \norm{\mathcal{T}(\vb) - \mathcal{T}(\vt)}_{\Ltwo{0, \widetilde{T}; \Ltwo{\Omega}}}^2 \leq q_2 \norm{\mathcal{T}_1(\vb) - \mathcal{T}_1(\vt)}_{\left(\Ltwo{0, \widetilde{T}; \Ltwo{\Gamma}}\right)^2}^2 \leq q_2 q_1 \norm{\vb - \vt}_{\Ltwo{0, \widetilde{T}; \Ltwo{\Omega}}}^2 ,
    \end{equation*}

    \noindent
    showing that $\mathcal{T}$ is a contraction since $q_2 q_1 \in [0, 1)$. We apply Banach's fixed-point theorem to get a unique fixed point $v \in \Ltwo{0, \widetilde{T}; \Ltwo{\Omega}}$. As previously discussed, this corresponds to a unique weak solution on $(0, \widetilde{T})$. 
    
    We now seek to extend this to a solution on $(0, T)$. Since all bounds for $u$, $u_b$ and $v$ were derived independently of $\widetilde{T}$, we can set new initial conditions
    %
    \begin{equation*}
        \vt_0 := v(\tfrac12 \widetilde{T}) \qquad \text{and} \qquad(\ut_0, \widetilde{u}_{b,0}) = (u(\tfrac12 \widetilde{T}), u_b(\tfrac12 \widetilde{T})) := \mathcal{T}_1(v)(\tfrac12 \widetilde{T})
    \end{equation*}
    
    \noindent
    and get the same contraction estimates. This gives another local solution $\vt$, this time defined on $(\frac12 \widetilde{T}, \frac32 \widetilde{T})$. Because of the uniqueness properties of both $v$ and $\vt$, their overlap on $(\frac12 \widetilde{T}, \widetilde{T})$ must coincide. This means that $\vt$ can be used to define an extension of $v$ to the interval $(0, \frac32 \widetilde{T})$.

    This process can be repeated, getting a solution on $(\frac{j}{2} \widetilde{T}, \frac{j + 2}{2} \widetilde{T})$ for $j = 2, \dots$ and using them to define extensions of $v$ until we cover the full time interval $(0, T)$. We can simultaneously construct the corresponding solutions to the ODEs by extending the local solutions using $(\ut, \ubt) := \mathcal{T}_1(\vt)$ at every step. For these solutions $u$, $u_b$ and $v$, the same non-negativity and boundedness properties hold as for all the local solutions, since these properties were proven independently of $\widetilde{T}$.
\end{proof}

\endgroup

% \begin{enumerate}
%     \item For every $v^* \in \mathcal{X}$, define $\mathcal{T}_{2} (v^*)$ as the solution $(u^*, u_b^*) \in \mathcal{Y}^2$ to the ODEs~\ref{temp:ode} with frozen reactions.
%     \item For every $(u^*, u_b^*) \in \mathcal{Y}^2$, define $\mathcal{T}_{1} (u^*, u_b^*)$ as the solution $v^* \in \mathcal{X}$ of the PDE~\ref{temp:pde} with frozen flux.
% \end{enumerate}

% We want to use the Picard-Lindel\"of theorem to prove existence and uniqueness of the ODEs. If we fix a solution of the PDEs $v$, we have a system

% \begin{equation}
%     \label{temp:ode}
%     \begin{aligned}
%         \partial_t u(t,x)    &= f(t, x, u(t,x), u_b(t,x)), \\
%         \partial_t u_b(t,x)  &= g(t, x, u(t,x), u_b(t,x)), \\
%     \end{aligned}
% \end{equation}

% \noindent
% where $f$ and $g$ both depend on the trace $v$. We need $f$ and $g$ to be Lipschitz continuous in $u$ and $u_b$ and continuous in time. The former follows from the definition of the model(???), but for the latter, we need continuity of $v$ on $\partial C$. In other words, for every $x \in \partial C$, we need the function $t \mapsto v \vert_{\partial C}(t, x)$ to be continuous. Then for every $x$, we can solve the system~\ref{temp:ode} for $u(\ \cdot\ , x)$ and $u_b(\ \cdot\ , x)$. What can we say about $u(t,\ \cdot\ )$ and $u_b(t,\ \cdot\ )$?

% For fixed $u$ and $u_b$, we have the following equations for $v$:

% \begin{equation}
%     \label{temp:pde}
%     \begin{aligned}
%         v_t(t, x) &= D\Delta v(t, x) && \text{on } (0, T) \times E,\\
%         D \pderiv[v]{\nu}(t, x) &= h(t, x, v(t, x)) && \text{on } (0, T) \times \partial C,\\
%         D \pderiv[v]{\nu}(t, x) &= 0 && \text{on } (0, T) \times \partial C,\\
%     \end{aligned}
% \end{equation}

% \noindent
% where $h$ depends on $u$ and $u_b$.

% \begin{equation*}
%     \begin{aligned}
%         \mathcal{X} &:= \left\{ v \in \Ltwo{0, T; \Omega} \cap C([0, T]; \Ltwo{\Omega}) \mid v \geq 0 \text{ and } v \vert_{\partial C}(\ \cdot\ , x) \in C([0, T]) \text{ for every } x \in \partial C \right\} \\
%         \mathcal{Y} &:= \left\{ u \in \Ltwo{0, T; H^1(\Omega)} \mid u \geq 0 \text{ and } u' \in \Ltwo{0, T; H^{-1}(\Omega)} \right\}
%     \end{aligned}
% \end{equation*}

% \noindent
% Questions:

% \begin{enumerate}
%     \item ODE is not Lipschitz, only locally?
%     \item Picard Lindel\"of for weak solutions?
%     \item Inward flow theorem for weak solutions?
%     \item When do we have time continuity of the trace?
%     \item System of ODEs $\rightarrow$ What can we say about $u(t,\ \cdot\ )$?
% \end{enumerate}